\documentclass[11pt,a4paper]{article}
\usepackage{fontspec,amsmath,amssymb,amsthm,microtype}
\setmainfont{lmroman10-regular.otf}[BoldFont=lmroman10-bold.otf,ItalicFont=lmroman10-italic.otf,BoldItalicFont=lmroman10-bolditalic.otf]
\usepackage[margin=25mm]{geometry}
\usepackage[unicode,hidelinks]{hyperref}
\hypersetup{pdfauthor={LangevinLean Project}}
\setlength{\emergencystretch}{3em}
\setlength{\parskip}{0.4em}
\setlength{\parindent}{0pt}
\allowdisplaybreaks[2]
\raggedbottom
\title{\texorpdfstring{Supporting planes and strong monotonicity}{Supporting planes and strong monotonicity}}
\author{LangevinLean Project}
\date{}
\begin{document}\maketitle
Draft mathematical exposition; these notes may contain mistakes.


Let \(H\) be a real inner-product space. We use the convention that \(f:H\to\mathbb R\) is \(m\)-strongly convex when \(v\mapsto f(v)-m\|v\|^2/2\) is convex. For a differentiable convex function, restriction to a line segment gives the supporting-plane inequality. Applying it to the quadratic-subtracted function gives the quadratic lower bound. Applying the lower bound to \((x,y)\) and \((y,x)\), and adding, gives strong monotonicity of the gradient.



The derivative \(Df(x)\) is a continuous linear functional. In a Hilbert space it is represented by \(\nabla f(x)\), so \(Df(x)[v]=\langle\nabla f(x),v\rangle\). These arguments require no second derivative or minimizer.


\section{\texorpdfstring{The supporting-plane inequality}{The supporting-plane inequality}}


Let \(H\) be a real inner-product space and let \(f:H\to\mathbb R\) be differentiable and convex on all of \(H\). For any \(x,y\in H\),

\[f(x)+Df(x)[y-x]\le f(y).\]

Restrict \(f\) to the affine line \(\ell(t)=x+t(y-x)\). Since affine composition preserves convexity, \(\varphi=f\circ\ell\) is convex on \(\mathbb R\), with \(\varphi(0)=f(x)\) and \(\varphi(1)=f(y)\).



The derivative of the line is its constant direction:

\[\ell'(0)=y-x.\]

By the chain rule, the derivative of the restricted function is

\[\varphi'(0)=Df(x)[y-x].\]

For a differentiable convex function, the derivative at the left endpoint is bounded by the secant slope to a later point. Applying this at \(0<1\) gives

\[Df(x)[y-x]\le\frac{\varphi(1)-\varphi(0)}{1-0}.\]

Evaluate the two endpoints and rearrange:

\[Df(x)[y-x]\le f(y)-f(x),\qquad f(x)+Df(x)[y-x]\le f(y).\]

The line argument uses only values on the segment joining the two points. For \(0<t\le1\), convexity gives \((\varphi(t)-\varphi(0))/t\le\varphi(1)-\varphi(0)\). Sending \(t\downarrow0\) identifies the left side with the directional derivative. The limit is the derivative at the segment parameter zero, and hence the derivative at the base point.


\section{\texorpdfstring{Strong monotonicity of the gradient}{Strong monotonicity of the gradient}}


Let \(H\) be a real Hilbert space. If \(f:H\to\mathbb R\) is differentiable and \(m\)-strongly convex, where \(m\in\mathbb R\), then for every \(x,y\in H\),

\[m\|x-y\|^2\le\langle\nabla f(x)-\nabla f(y),x-y\rangle.\]

Set

\[g(v)=f(v)-\frac m2\|v\|^2.\]

The definition of strong convexity says precisely that \(g\) is convex.



The squared norm has derivative \(D\|v\|^2[u]=2\langle v,u\rangle\). Consequently \(g\) is differentiable and

\[Dg(v)[u]=Df(v)[u]-m\langle v,u\rangle.\]

Apply the supporting-plane inequality in both directions, and substitute the derivative just computed:

\[g(x)+Dg(x)[y-x]\le g(y),\qquad g(y)+Dg(y)[x-y]\le g(x).\]

Express the derivative of \(f\) using its gradient. The two inequalities become

\[\begin{aligned}f(x)-\tfrac m2\|x\|^2+\langle\nabla f(x),y-x\rangle-m\langle x,y-x\rangle&\le f(y)-\tfrac m2\|y\|^2,\\ f(y)-\tfrac m2\|y\|^2+\langle\nabla f(y),x-y\rangle-m\langle y,x-y\rangle&\le f(x)-\tfrac m2\|x\|^2.\end{aligned}\]

Add these inequalities. The function values and quadratic values cancel. Using \(\langle x-y,x-y\rangle=\|x-y\|^2\), what remains is

\[0\le\langle\nabla f(x)-\nabla f(y),x-y\rangle-m\|x-y\|^2.\]

This is the required monotonicity inequality.



Reversing the pair changes \(y-x\) into \(-(x-y)\). Therefore adding the two supporting inequalities makes the function values cancel and leaves a gradient difference paired with a displacement of the same orientation. The resulting nonnegative pairing gives monotonicity of \(\nabla f\), rather than the opposite inequality. No Hessian is introduced.


\section{\texorpdfstring{The quadratic lower model}{The quadratic lower model}}


Let \(H\) be a real Hilbert space, \(f:H\to\mathbb R\) differentiable, and \(m\in\mathbb R\). If \(f\) is \(m\)-strongly convex on all of \(H\), then for every \(x,y\in H\),

\[f(x)+\langle\nabla f(x),y-x\rangle+\frac m2\|y-x\|^2\le f(y).\]

Set \(g(v)=f(v)-\frac m2\|v\|^2\). Strong convexity makes \(g\) convex.



Differentiate this shifted potential:

\[Dg(v)[u]=Df(v)[u]-m\langle v,u\rangle.\]

In particular, \(g\) is differentiable on \(H\).



Its supporting-plane inequality at \(x\), together with \(Df(x)[u]=\langle\nabla f(x),u\rangle\), yields

\[f(x)-\frac m2\|x\|^2+\langle\nabla f(x),y-x\rangle-m\langle x,y-x\rangle\le f(y)-\frac m2\|y\|^2.\]

Collect the quadratic correction by the identity

\[\|y\|^2-\|x\|^2-2\langle x,y-x\rangle=\|y-x\|^2.\]

Moving terms gives the desired lower model.



The correction is the squared displacement, not the difference of the two squared norms. The cross term from differentiating the subtracted quadratic is exactly what converts that difference into \(\|y-x\|^2\). When \(\nabla f(x)=0\), this inequality specializes to \(f(y)-f(x)\ge m\|y-x\|^2/2\), which gives a Gaussian upper bound on the density proportional to the exponential of the negative potential.


\end{document}
